\documentclass[11pt]{article}
\usepackage[a4paper,margin=1in]{geometry}
\usepackage[T1]{fontenc}
\usepackage{lmodern}
\usepackage{microtype}
\usepackage{amsmath,amssymb,amsthm,mathtools}
\usepackage{booktabs}
\usepackage{enumitem}
\usepackage{xcolor}
\usepackage[numbers]{natbib}
\usepackage[colorlinks=true,allcolors=blue]{hyperref}

\newtheorem{theorem}{Theorem}
\newtheorem{corollary}[theorem]{Corollary}
\newtheorem{proposition}[theorem]{Proposition}
\newcommand{\C}{\mathbb C}
\newcommand{\J}{\mathcal J}

\title{Anchored Shell-Prefix Availability Obstruction}
\author{Bin Wang}
\date{July 2026}

\begin{document}
\maketitle

\begin{abstract}
The shell-prefix selector of RH-238 makes a finite cloud-extracted trace jet
small relative to zero.  The deterministic numerator dictionary of RH-243
instead prescribes a nonzero coefficient anchor.  We first prove that the two
equal-tolerance jet balls are disjoint throughout the archived noise range.
We then scan every shell-complete prefix in the frozen RH-222 candidate
windows against the anchor: none of 543 prefixes passes at any of 32
endpoints.  This is a finite obstruction for one candidate class and one
tolerance rule, not a nonexistence theorem for alternative clouds.
\end{abstract}

\section{Two selection problems}

For jets $x=(x_2,\ldots,x_N)$ and $y=(y_2,\ldots,y_N)$ define
\[
 d_N(x,y)=\sum_{n=2}^{N}\frac{|x_n-y_n|}{n},
 \qquad \J_N(x)=d_N(x,0).
\]
This is the radius-one logarithmic jet metric used in RH-237 and RH-238
\citep{WangRH238}.  Let $a=(a_2,\ldots,a_{12})$ be the Hardy-scaled
deterministic numerator anchor of RH-243 \citep{WangRH243}.  Its norm is
\begin{equation}\label{eq:anchor-norm}
 \J_{12}(a)=0.49450543569144195.
\end{equation}

The RH-238 condition is $\J_{12}(\tau)\le\varepsilon_\sigma$.  The anchored
condition required for coefficient identification is instead
$d_{12}(\tau,a)\le\varepsilon_\sigma$.  The following elementary separation
is decisive before any candidate scan.

\begin{theorem}[disjoint target balls]\label{thm:balls}
For every normed jet space and every nonzero anchor $a$, the closed balls
$B(0,\varepsilon)$ and $B(a,\varepsilon)$ are disjoint if
$\|a\|>2\varepsilon$.  More quantitatively, if $\|x\|\le\varepsilon$, then
\[
 \|x-a\|\ge \|a\|-\varepsilon.
\]
\end{theorem}

\begin{proof}
The reverse triangle inequality gives
$\|x-a\|\ge|\|a\|-\|x\||\ge\|a\|-\varepsilon$.  If $x$ belonged to both
equal-radius balls, the ordinary triangle inequality would give
$\|a\|\le2\varepsilon$, a contradiction.
\end{proof}

\begin{corollary}[archived tolerance separation]\label{cor:archive}
For $\varepsilon_\sigma=\sigma$ and every archived
$0.00125\le\sigma\le0.04$, a jet satisfying the RH-238 zero-target condition
cannot satisfy the RH-243 anchored condition at the same tolerance.  The
smallest strict separation margin $\J_{12}(a)-2\sigma$ is
$0.41450543569144194$.
\end{corollary}

Thus the earlier 32/32 zero-target success is not evidence for anchored
availability.  This conclusion is exact for the displayed finite metric and
does not depend on floating root locations.

\section{Frozen prefix scan}

We next ask the different question whether some other prefix in the same
finite candidate class meets the anchor.  At each noise/side endpoint let
$A_\sigma$ be the archived scaled matrix, $p_\sigma$ and $e_\sigma$ its Perron
and parity roots, and
\[
 S_{\sigma,1},\ldots,S_{\sigma,m_\sigma}
\]
the conjugation-complete shells obtained from the ordered RH-222 candidate
window \citep{WangRH222}.  For prefix $k$ put
\[
 \tau_n^{(k)}
 =\operatorname{tr}A_\sigma^n-p_\sigma^n-e_\sigma^n
  -\sum_{s\in S_{\sigma,1}\cup\cdots\cup S_{\sigma,k}}s^n,
 \qquad 2\le n\le12.
\]
We retain RH-238's minimum rank four and scan every eligible prefix, but now
test
\begin{equation}\label{eq:test}
 d_{12}(\tau^{(k)},a)\le\sigma.
\end{equation}
The full traces come from RH-236 \citep{WangRH236}; the anchor is read from
the independent RH-243 archive.

\begin{proposition}[finite anchored prefix obstruction]\label{prop:scan}
Across the 32 frozen endpoints, all 543 eligible shell-complete prefixes fail
\eqref{eq:test}.  Endpointwise best distances range from
$0.39723767197524446$ to $0.48457639371229216$.  Their ranks range from 4 to
18.  The minimum and maximum ratios of best distance to tolerance are
$10.479757156015458$ and $343.82932788167034$.
\end{proposition}

\begin{proof}
The accompanying deterministic audit reconstructs the conjugate shells,
forms every eligible cumulative prefix, evaluates the exact finite spectral
power sums, and applies the displayed weighted metric.  The archived JSON
records every endpointwise optimum and count.  Unit tests independently
check the trace subtraction, metric, target-ball theorem, and aggregate
failure ledger.
\end{proof}

\begin{center}
\small
\begin{tabular}{lr}
\toprule
quantity & archived value\\
\midrule
endpoints & 32\\
eligible prefixes & 543\\
anchored passes & 0\\
best-distance range & $0.397238$--$0.484576$\\
best-rank range & 4--18\\
minimum best-distance/tolerance & $10.479757$\\
\bottomrule
\end{tabular}
\end{center}

\section{Scope and route consequence}

Proposition~\ref{prop:scan} is exhaustive only inside the frozen RH-222
candidate windows, their radial shell ordering, the shell-complete prefix
class, orders 2--12, radius one, minimum rank four, and the rule
$\varepsilon_\sigma=\sigma$.  The Arnoldi roots are floating approximations,
not interval enclosures.  The result does not exclude non-prefix subsets,
weighted cancellations, expanded candidate windows, different tolerances, or
a continuum small-noise selector.  It therefore cannot be promoted to an
asymptotic cloud nonexistence theorem.

The useful route consequence is narrower: zero-target compression and
deterministic numerator identification must remain separate obligations.  A
new anchored selector is needed.  Independently, the all-order envelope
should be pursued through cancellation-preserving loop grouping rather than
absolute physical and atomic sectors.  Gate A remains open, and Gates B--E
are untouched.  No Hilbert--P\'olya operator, zeta-divisor identification,
Riemann-zero identification, or Riemann-hypothesis implication is claimed.

\bibliographystyle{plainnat}
\bibliography{references}
\end{document}
